%%%%%%%%%%%%%%%%%%%%%%%%%%%%%%%%%%%%%%%%%
% Medium Length Professional CV
% LaTeX Template
% Version 2.0 (8/5/13)
%
% This template has been downloaded from:
% http://www.LaTeXTemplates.com
%
% Original author:
% Rishi Shah 
%
% Important note:
% This template requires the resume.cls file to be in the same directory as the
% .tex file. The resume.cls file provides the resume style used for structuring the
% document.
%
%%%%%%%%%%%%%%%%%%%%%%%%%%%%%%%%%%%%%%%%%

%----------------------------------------------------------------------------------------
%	PACKAGES AND OTHER DOCUMENT CONFIGURATIONS
%----------------------------------------------------------------------------------------

\documentclass{resume} % Use the custom resume.cls style

\usepackage[margin=1.0in]{geometry} % Document margins
\newcommand{\tab}[1]{\hspace{.2667\textwidth}\rlap{#1}}
\newcommand{\itab}[1]{\hspace{0em}\rlap{#1}}
\name{Kentaro Sugimoto} % Your name
\address{Nishisugamo, Toshima-ku, Tokyo, Japan} % Your address
\address{\#304 International Building A, 2-31-7 \\ 170-0001} % Your secondary addess (optional)
\address{(+81)90-1717-9903 \\ tarotene@iis.u-tokyo.ac.jp} % Your phone number and email

\usepackage{color}
\usepackage{url}
\usepackage{hyperref}

\begin{document}

\begin{rSection}{Skills}

\begin{tabular}{ @{} >{\bfseries}l @{\hspace{6ex}} l }
Language \ & Japanese (Native), English (Advanced)\\
Scientific Skills \ & Statistical Analysis, Basic Mathematics\\
Programming Language \ & C/C++, Fortran, Python, Unix Shell
\end{tabular}

\end{rSection}

\begin{rSection}{Work Experience}

\begin{rSubsection}{Cluster for Pioneering Research, RIKEN}{April 2018 - Present}{Junior Research Associate (Host Researcher: Seiji Yunoki)}{}
\item Conducting a research on \textit{spin-state manipulation in magnetic materials}.
\end{rSubsection}

\begin{rSubsection}{MDR Inc. (Machine Learning \& Dynamics Research)}{March 2019}{Student Internship}{}
\item Developed extensions for \textit{Blueqat} (the company's quantum-computing software).
\end{rSubsection}

\begin{rSubsection}{College of Arts and Sciences, Univ. Tokyo}{Apr. 2017 - Sep. 2017}{Teaching Assistant for the cource Information}{}
\item Assisted students to use unix shell environments and do their exercises.
\end{rSubsection}

\end{rSection}

\begin{rSection}{Education}

{\bf University of Tokyo} \hfill {\em April 2018 - Present} 
\\ Doctor of Science, Department of Physics (Expected completion)\\
\\{\bf University of Tokyo} \hfill {\em April 2016 - March 2018} 
\\ Master of Science, Department of Physics
\\ Thesis: \textit{Effects of Boundary Conditions on Magnetic Friction}\\
\\{\bf Osaka Prefecture University} \hfill {\em April 2011 - March 2016} 
\\ Bachelor of Science, Department of Physical Science
\\ Thesis: \textit{Formulation of a Linearized Dispersion Equation and its Applications}

\end{rSection}

\newpage

\begin{rSection}{Completed Courseworks}

\begin{rSubsection}{Innovative Professional Development Program}{August 2019 - March 2020}{Professional development Consortium for Computational Materials Scientists}{}
\item Acquired knowledge and techniques of scientific/high-performance computing.
\item Deeply understand corporate needs and international research trends.
\end{rSubsection}

\begin{rSubsection}{Future Faculty Program (FFP)}{April 2019 - September 2019}{Graduate School of Arts and Sciences, Univ. Tokyo}{}
\item Learned the current situation of higher education, lesson design, and active learning techniques.
\end{rSubsection}

\end{rSection}

\begin{rSection}{Ongoing Courseworks}

\begin{rSubsection}{Computational Science Alliance}{April 2018 - Present}{}{}
\item Learned various parallel-programming and tuning techniques.
\item Learned finite element method, singular value decomposition, and their mathematical foundations.
\end{rSubsection}

\begin{rSubsection}{Data Scientist School (DSS)}{April 2019 - Present}{Education and Research Program for Domain-Knowledge Creation, Univ. Tokyo}{}
\item Analyzed data given by corporations in a group-work style.
\item Learned problem-discovering, suggesting the way of analysis, actual analysis, and the feedback.
\end{rSubsection}

\end{rSection}

\begin{rSection}{Publications and Talks}{}{}{}
    \begin{rSubsection}{Journal Article (Reviewed)}{}{}{}
        \item K. Sugimoto, Phys. Rev. E \textbf{99}, 052103 (2019)\\
        $[$\url{https://journals.aps.org/pre/abstract/10.1103/PhysRevE.99.052103}$]$
    \end{rSubsection}

    \begin{rSubsection}{International Conference (Reviewed)}{}{}{}
        \item K. Sugimoto, M. Matsuo, StatPhys27 (2019) (Poster)\\
        $[$\url{https://statphys27.df.uba.ar/registration/index.php/SP27/MainConference/paper/view/498}$]$
    \end{rSubsection}

    \begin{rSubsection}{Domestic Conference}{}{}{}
        \item K. Sugimoto, N. Hatano, JPSJ the 72th Annual Meeting (2017)\\
        $[$\url{https://doi.org/10.11316/jpsgaiyo.72.1.0_2909}$]$
        \item K. Sugimoto, N. Hatano, JPSJ 2017 Autumn Meeting (2017)\\
        $[$\url{https://doi.org/10.11316/jpsgaiyo.72.2.0_2593}$]$
        \item K. Sugimoto, N. Hatano, JPSJ the 73th Annual Meeting (2018)\\
        $[$\url{https://doi.org/10.11316/jpsgaiyo.73.1.0_2757}$]$
        \item K. Sugimoto, M. Matsuo, JPSJ the 74th Annual Meeting (2019)
    \end{rSubsection}
\end{rSection}

\begin{rSection}{Reference}{}{}{}
Available upon request.
\end{rSection}

\end{document}
